\section{\emph{Celestial Hierarchy} Citations in the \emph{Summa theologiae}}
\label{app:ch-concordance}

This concordance maps all fifteen chapters of Pseudo-Dionysius's \emph{De caelesti hierarchia} in exact order, giving the chapter title from John Parker's 1899 translation, the Migne \emph{Patrologia Graeca}~3 column span of Dionysius's own text (excluding interleaved Corderius notes and Pachymeres paraphrase), the Parker page and section bounds, and the exact \emph{Summa theologiae} articles where Aquinas receives, adapts, or qualifies each chapter, in exact agreement with Section~\ref{sec:celestial-hierarchy}.

{\footnotesize
\setlength{\tabcolsep}{2.5pt}
\renewcommand{\arraystretch}{1.15}
\begin{longtable}{>{\bfseries\raggedright\arraybackslash}p{0.04\linewidth}>{\raggedright\arraybackslash}p{0.24\linewidth}>{\raggedright\arraybackslash}p{0.10\linewidth}>{\raggedright\arraybackslash}p{0.10\linewidth}>{\raggedright\arraybackslash}p{0.21\linewidth}>{\raggedright\arraybackslash}p{0.23\linewidth}}
\toprule
\textbf{Ch.} & \textbf{Parker (1899) Chapter Title} & \textbf{PG~3 Cols.} & \textbf{Parker pp.\ / \S\S} & \textbf{\emph{ST} Reception Loci} & \textbf{How Aquinas Receives or Qualifies It} \\
\midrule
\endfirsthead
\toprule
\textbf{Ch.} & \textbf{Parker (1899) Chapter Title} & \textbf{PG~3 Cols.} & \textbf{Parker pp.\ / \S\S} & \textbf{\emph{ST} Reception Loci} & \textbf{How Aquinas Receives or Qualifies It} \\
\midrule
\endhead
\bottomrule
\endfoot
I & That every divine illumination, whilst going forth lovingly to the objects of its forethought under various forms, remains simplex, and unifies the illuminated & 119--124 (120B--124A) & pp.~1--4; \S\S~I--III & \ST{I, q.~1, a.~9; q.~108, a.~1; q.~111, a.~1} & Quotes \CH{1.2} on sacred veils for scriptural metaphor; grounds unity of angelic hierarchy and descent of divine light to men. \\
\addlinespace[2pt]
II & That Divine and Heavenly things are appropriately revealed, even through dissimilar symbols & 136--146 & pp.~4--13; \S\S~I--V & \ST{I, q.~1, a.~9, ad~3} & Adopts Dionysius's three reasons for dissimilar corporeal metaphors as protective, apophatic pedagogy in Scripture. \\
\addlinespace[2pt]
III & What is Hierarchy? and what the use of Hierarchy? & 164--168 & pp.~13--16; \S\S~I--III & \ST{I, q.~106, aa.~1, 3; q.~108, aa.~1, 2, 4} & Receives definition of hierarchy (sacred order, knowledge, activity) and threefold purification, illumination, perfection; clarifies no impurity exists in blessed angels. \\
\addlinespace[2pt]
IV & What is meant by the appellation ``Angels?'' & 177--184 & pp.~16--21; \S\S~I--IV & \ST{I, q.~57, aa.~1, 5; q.~106, a.~3; q.~111, a.~1; q.~113, a.~2; III, q.~30, a.~2} & Adopts hierarchical law of mediated revelation (including Gabriel's Annunciation and Mosaic Law); applies ``Angel of Great Counsel'' to Christ's mission. \\
\addlinespace[2pt]
V & For what reason all the Heavenly Beings are called, in common, Angels & 196 & pp.~21--22; undivided & \ST{I, q.~108, a.~5; q.~113, a.~3} & Uses common-vs.-proper naming rule: higher orders possess all lower perfections eminently while lowest order takes the common name ``Angels.'' \\
\addlinespace[2pt]
VI & Which is the first Order of the Heavenly Beings? which the middle? and which the last? & 200--201 & pp.~23--24; undivided & \ST{I, q.~56, a.~1; q.~63, a.~7; q.~108, aa.~1--3} & Receives the three-triad structure of nine choirs; compares old and new Latin translations on angelic self-knowledge (q.~56, a.~1). \\
\addlinespace[2pt]
VII & Concerning the Seraphim and Cherubim and Thrones, and concerning their first Hierarchy & 205--212 & pp.~24--31; \S\S~I--IV & \ST{I, q.~57, a.~5; q.~63, a.~7; q.~106, aa.~1--2, 4; q.~107, a.~2; q.~108, aa.~1, 4--6; q.~112, a.~4} & Central source for the First Triad's etymologies, direct vision of God, angelic questioning of the Ascension (Is~63:1), and downward illumination. \\
\addlinespace[2pt]
VIII & Concerning Lordships and Powers and Authorities, and concerning their middle Hierarchy & 237--242 & pp.~31--35; \S\S~I--II & \ST{I, q.~106, a.~1; q.~108, aa.~5--6; q.~112, a.~4} & Adopts Middle Triad order (Dominations, Virtues, Powers) and defines Dominations as directing without external mission while reconciling Gregory. \\
\addlinespace[2pt]
IX & Concerning the Principalities, Archangels, and Angels, and concerning their last Hierarchy & 257--264 & pp.~35--40; \S\S~I--IV & \ST{I, q.~108, aa.~5--6; q.~113, a.~3} & Grounds national governance by Principalities (Deut~32:8 LXX; Dan~10) and individual guardianship by the lowest order (Angels). \\
\addlinespace[2pt]
X & A Repetition and Summary of the Angelic discipline & 272--274 & pp.~41--42; \S\S~I--III & \ST{I, q.~50, a.~4; q.~108, a.~3; q.~113, aa.~2--3} & Invokes internal first/middle/last gradation within each order alongside species-distinctness (q.~50, a.~4) and ordered guardianship. \\
\addlinespace[2pt]
XI & For what reason all the Heavenly Beings, in common, are called Heavenly Powers & 284--285 & pp.~42--44; \S\S~I--II & \ST{I, q.~54, a.~3} & Draws the triad \emph{essence, power, operation} (\gk{ousia, dynamis, energeia}) directly into his proof that an angel's power differs from its essence. \\
\addlinespace[2pt]
XII & Why the Hierarchs amongst men are called Angels & 292--293 & pp.~44--46; \S\S~I--III & \ST{I, q.~55, a.~3; q.~106, a.~4} & Cites the hierarchical principle of higher participation in universal light to establish fewer, more universal species in higher angels. \\
\addlinespace[2pt]
XIII & For what reason the Prophet Isaiah is said to have been purified by the Seraphim & 300--308 & pp.~46--53; \S\S~I--IV & \ST{I, q.~112, a.~2; q.~113, a.~2} & Key locus on whether assisting Seraphim are sent externally: Aquinas follows Dionysius's principal solution (a lower angel acting in the Seraph's power and name) while integrating Gregory's \emph{Hom.}~34, and cites it on mediated guardianship. \\
\addlinespace[2pt]
XIV & What the traditional number of the Angels signifies & 321--322 & pp.~53--54; undivided & \ST{I, q.~50, a.~3; q.~112, a.~4} & Quotes Dionysius's reading of Dan~7:10 directly to establish that angelic multitude vastly surpasses all material number while remaining finite to God. \\
\addlinespace[2pt]
XV & What are the morphic likenesses of the Angelic Powers? (fire, human form, wings, wheels, rivers, joy) & 325--340 & pp.~54--66; \S\S~I--IX & \ST{I, q.~51, a.~3, ad~4; q.~106, aa.~1, 4; q.~107, a.~5} & Quotes the division and multiplication of uniform light in illumination (q.~106, aa.~1, 4), selective communication (q.~107, a.~5), and symbolic readings of bodily limbs attributed to angels in Scripture (q.~51, a.~3, ad~4). \\
\end{longtable}
}
